\writeSectionFrame{\presentationTitle}{Das Iterator-Muster}
\begin{frame}{Steckbrief}
	\begin{itemize}
 		\item Deutscher Name: Iterator
 		\item Alternativer Name: Cursor
 		\item Englischer Name: Iterator
 		\item Gruppe: Verhaltensmuster
 	\end{itemize}
\end{frame}

\begin{frame}{Iterator}
	\begin{block}{Zweck}
 		Biete eine Möglichkeit, um auf die Elemente eines zusammengesetzten Objekts sequentiell zugreifen zu können, ohne die zugrundeliegende Repräsentation offenzulegen.
 	\end{block}
\end{frame}

\begin{frame}{Allgemeines Klassendiagramm}
	\begin{center}
 		\includegraphics[scale=0.45]{\BILDERPFAD/iterator1.png}
 	\end{center}
\end{frame}
\note{
	Die Bezeichnungen können am Anfang etwas verwirren, aber sie sind halt sehr
	generisch. Ein Aggregat ist eine Zusammenfassung mehrere Werte. Das kann eine
	Liste sein, aber auch bspw. ein Datenbankcursor. Alle diese iterierbaren Klassen
	implementieren ein Interface (natürlich auch über mehrere Stufen möglich).  Dieses 
	Interface hat nur eine einzige Methode, die einen Iterator erzeugt. Die gemeinsame 
	Schnittstelle der Iteratoren ist das Interface Iterator. erzeugeIterator() gibt 
	einen konkreten Iterator zurück, der die Methoden zum Durchlaufen durch die Menge
	implementiert. Dazu muss der konkrete Iterator das konkrete Aggregat kennen. Der
	Klient lässt sich vom konkreten Aggregat einen konkreten Iterator erzeugen und
	benutzt diesen, um die einzelnen Elemente zu durchlaufen. Damit weiß er von der
	tatsächlichen Logik zum Durchlaufen der Elemente nichts. 	
}

\begin{frame}{Klassendiagramm - Iterator im JDK}
	\begin{center}
 		\includegraphics[scale=0.35]{\BILDERPFAD/iterator2.png}
 	\end{center}
\end{frame}

\note{
	Die Implementierung des Iterator-Musters im JDK entspricht nicht ganz dem vorgestellten
	Muster. Die Iterator-Schnittstelle unterstützt die Methode remove(). Die Standardimplementierung
	wirft eine UnsupportedOperationException, wenn sie vom konkreten Iterator nicht geändert wird.
	Der Iterator kann nicht auf das erste Element gesetzt werden, was aber kein Problem ist: Man
	kann sich einfach ein neues Iteratorobjekt erzeugen lassen. Auch kann man das aktuelle Element
	nicht mit einer Methode anfordern. Aber auch das ist kein Problem, weil die Methode next() das
	aktuelle Element zurückgibt. Man muss es sich nur selbst merken. Die Methode forEachRemaining()
	erhält eine Operation, die für die restlichen Elemente ausgeführt wird. Die konkrete Iterator-Klassen
	Itr ist privat, weil der Iterator ja nur für die ArrayList arbeitet und eng mit ihr verbunden ist. 
	So kann der Iterator für die ArrayList optimiert werden. Die zweite
	Schnittstelle ListIterator stellt weitere Funktionen zur Verfügung wie z.B. das Rückwärtslaufen.
}

\begin{frame}[fragile]{Noch mehr JDK}
	\begin{exampleblock}{StandardIteratorBeispiel.java}
 		\begin{lstlisting}
Iterator<String> zauberspruchIterator = 
        zauberspruch.iterator();
while(zauberspruchIterator.hasNext()) {
    String myString = zauberspruchIterator.next();
    System.out.print(myString);
    zauberspruch.add("Und dreimal schwarzer Kater");
}
        \end{lstlisting}
 	\end{exampleblock}
\end{frame}

\note{
	Nun stellt sich die Frage, warum der Iterator Methoden anbietet zum Hinzufügen und Löschen von Elementen. Kann das
	nicht das Aggregat selbst machen? Dafür muss genau klar sein, was ein Iterator tut. Er durchläuft eine Menge von 
	Elementen und merkt sich die aktuelle Position. Er muss sich auf die zugrundeliegende Datenquelle verlassen
	können. Wenn die jetzt einfach verändert wird, gerät der Iterator in Schwierigkeiten. Deswegen wird in diesem
	Beispiel eine ConcurrentModificationException geworfen und Änderungen dürfen nur über den Iterator durchgeführt
	werden. Dann wird der Iterator nämlich informiert und er kann die Änderung berücksichtigen.
}

\frameWithImageOnly{\BILDERPFAD/iterator3.jpg}
\note{
	Wir möchte gerne einen MP3-Player simulieren. Einen Iterator zu bauen, der einfach nur durch die Lieder läuft
	ist eher langweilig, weil List diesen ja schon liefert. Interessanter ist ein Iterator, der die Lieder in
	zufälliger Reihenfolge abspielt. Diese Zufallslogik soll im Iterator gekapselt werden.
}

\begin{frame}{Klassendiagramm - MP3Player}
	\begin{center}
 		\includegraphics[scale=0.3]{\BILDERPFAD/iterator4.png}
 	\end{center}
\end{frame}

\note{
	Der Iterator wird in der ersten Version mit E typisiert, damit wir auch mehr als nur MP3s verwenden können.
	Erst der MP3Player weiß tatsächlich, dass wir mit MP3s arbeiten.
}

\begin{frame}[fragile]{Iterator}
	\begin{exampleblock}{Iterator}
 		\begin{lstlisting}
public interface Iterator<E> {
	void start();
	void next();
	boolean finished();
	E currentElement();
}
 		\end{lstlisting}
 	\end{exampleblock}
\end{frame}

\begin{frame}[fragile]{Konkreter Iterator}
	\begin{exampleblock}{\javapackageEins{RandomIterator}}
 		\begin{lstlisting}
public class RandomIterator<E> implements Iterator<E> {
	private List<E> originalList;
	private List<E> randomList;
	private int currentPosition;
	
	public RandomIterator(List<E> originalList) {
		this.originalList = originalList;
		start();
	}
 		\end{lstlisting}
 	\end{exampleblock}
\end{frame}

\begin{frame}[fragile]{Konkreter Iterator}
	\begin{exampleblock}{\javapackageEins{RandomIterator}}
 		\begin{lstlisting}
@Override
public void next() {
    if (!finished()) {
        currentPosition++;
    }
}

@Override
public boolean finished() {
    if (currentPosition >= (randomList.size() - 1)) {
        return true;
    } else {
        return false;
    }
}
 		\end{lstlisting}
 	\end{exampleblock}
\end{frame}

\begin{frame}[fragile]{Konkreter Iterator}
	\begin{exampleblock}{\javapackageEins{RandomIterator}}
 		\begin{lstlisting}
@Override
public E currentElement() {
    if (currentPosition == -1) {
        return null;
    } else {
        return randomList.get(currentPosition);
    }
} 		    
 		\end{lstlisting}
 	\end{exampleblock}
\end{frame}

\begin{frame}[fragile]{Aggregat}
	\begin{exampleblock}{\javapackageEins{Playable}}
 		\begin{lstlisting}
public interface Playable<E> {
	Iterator<E> createIterator();
}
 		\end{lstlisting}
 	\end{exampleblock}
\end{frame}

\begin{frame}[fragile]{Konkretes Aggregat}
	\begin{exampleblock}{\javapackageEins{MusicLibrary}}
 		\begin{lstlisting}
public class MusicLibrary<E> implements Playable<E> {
	private List<E> music;
	
	public MusicLibrary() {
		music = new ArrayList<>();
	}

	public void addSong(E song) {
		music.add(song);
	}
 		\end{lstlisting}
 	\end{exampleblock}
\end{frame}

\begin{frame}[fragile]{Konkretes Aggregat}
	\begin{exampleblock}{\javapackageEins{MusicLibrary}}
 		\begin{lstlisting}
@Override
public Iterator<E> createIterator() {
    return new RandomIterator<E>(music);
}

public void play(E song) {
    System.out.println("Im Moment laeuft: " + song.toString());
}
 		\end{lstlisting}
 	\end{exampleblock}
\end{frame}

\begin{frame}{Was ist mit anderen Formaten?}
	\begin{itemize}
 		\item MP3
 		\item CD
 		\item WMA
 		\item OGG
 		\item Ganz neue Formate
 	\end{itemize}
\end{frame}

\begin{frame}{Klassendiagramm - Verbesserung}
	\begin{center}
 		\includegraphics[scale=0.3]{\BILDERPFAD/iterator5.png}
 	\end{center}
\end{frame}

\note{
	Hier können wir auf die Stärke von Interfaces zurückgreifen, indem wir der Klasse MP3 ein Interface Music
	verpassen. Wenn wir nun ein anderes Format einbauen möchten, bilden wir einfach eine neue implementierende
	Klasse. Mit einem anderen Interface wie z.B. Medium könnte man es noch generischer halten und z.B. auch
	MP3s abspielen.
}

\frameWithImageOnly{\BILDERPFAD/iterator6.jpg}
\note{
	Wir möchten jetzt einen zweiten Iterator einbauen. Dieser bekommt den Namen eines Interpreten und dann
	möchten wir nur die Musik von diesem Interpreten hören.
}

\begin{frame}{Klassendiagramm - Neuer Iterator}
	\begin{center}
 		\includegraphics[scale=0.3]{\BILDERPFAD/iterator7.png}
 	\end{center}
\end{frame}

\note{
	Um einen neuen Iterator einzubauen, müssen wir tatsächlich nur einen neuen konkreten Iterator einführen. 
	Mehr Änderungen sind nicht notwendig. 
}

\begin{frame}[fragile]{Neuer konkreter Iterator}
	\begin{exampleblock}{\javapackageZwei{ArtistIterator}}
 		\begin{lstlisting}
public class ArtistIterator implements Iterator<Music> {
	private List<Music> originalList;
	private int currentPosition;
	private String artist;
	
	public ArtistIterator(List<Music> originalList, 
			String artist) {
		this.originalList = originalList;
		this.artist = artist;
		start();
	}

	@Override
	public void start() {
		currentPosition = -1;
	}
 		\end{lstlisting}
 	\end{exampleblock}
\end{frame}

\begin{frame}[fragile]{Neuer konkreter Iterator}
	\begin{exampleblock}{\javapackageZwei{ArtistIterator}}
 		\begin{lstlisting}
@Override
public void next() {
    if (!finished()) {
        currentPosition++;
    }
}

@Override
public boolean finished() {
    if (currentPosition >= (originalList.size() - 1)) {
        return true;
    } else {
        return false;
    }
}
 		\end{lstlisting}
 	\end{exampleblock}
\end{frame}

\begin{frame}[fragile]{Neuer konkreter Iterator}
	\begin{exampleblock}{\javapackageZwei{ArtistIterator}}
 		\begin{lstlisting}
@Override
public Music currentElement() {
    if (currentPosition == -1) {
        return null;
    } else {
        Music music = originalList.get(currentPosition);
        
        if (music.getArtist().equals(artist)) {
            return originalList.get(currentPosition);
        }
    }
    
    return null;
}
 		\end{lstlisting}
 	\end{exampleblock}
\end{frame}


\begin{frame}{Iterator-Arten}
	\begin{itemize}
 		\item extern
 		\item intern
 	\end{itemize}
\end{frame}

\note{
	Wir haben in unserem Beispiel einen externen Iterator entwickelt, d.h. der Klient steuer die Iterierung. Er ruft also
	die next()-Methode selbst auf. Beim	internen Iterator wird sie vom Iterator selbst gesteuert. 
}

\begin{frame}[fragile]{Interner Iterator}
	\begin{exampleblock}{\javapackageDrei{Iterator}}
 		\begin{lstlisting}
public interface Iterator<Music> {
	void start();
	void next();
	boolean finished();
	Music currentElement();
	void iterate(Consumer<Music> action);
} 		    
 		\end{lstlisting}
 	\end{exampleblock}
\end{frame}

\note{
	Mit Java 8 ist ein interner Iterator recht einfach zu implementieren.
}

\begin{frame}[fragile]{Interner Iterator}
	\begin{exampleblock}{\javapackageDrei{RandomIterator}}
 		\begin{lstlisting}
public class RandomIterator implements Iterator<Music> {
	private List<Music> originalList;
	private List<Music> randomList;
	private int currentPosition;
	
	public RandomIterator(List<Music> originalList) {
		this.originalList = originalList;
		start();
	}

	@Override
	public void start() {
		randomList = new ArrayList<>(originalList);
		Collections.shuffle(randomList);
		currentPosition = -1;
	}
 		\end{lstlisting}
 	\end{exampleblock}
\end{frame}

\begin{frame}[fragile]{Interner Iterator}
	\begin{exampleblock}{\javapackageDrei{RandomIterator}}
 		\begin{lstlisting}
@Override
public void next() {
    if (!finished()) {
        currentPosition++;
    }
}

@Override
public boolean finished() {
    if (currentPosition >= (randomList.size() - 1)) {
        return true;
    } else {
        return false;
    }
}
 		\end{lstlisting}
 	\end{exampleblock}
\end{frame}

\begin{frame}[fragile]{Interner Iterator}
	\begin{exampleblock}{\javapackageDrei{RandomIterator}}
 		\begin{lstlisting}
@Override
public Music currentElement() {
    if (currentPosition == -1) {
        return null;
    } else {
        return randomList.get(currentPosition);
    }
}

@Override
public void iterate(Consumer<Music> action) {
    while (!finished()) {
        next();
        action.accept(currentElement());
    }
}
 		\end{lstlisting}
 	\end{exampleblock}
\end{frame}

\begin{frame}[fragile]{Klient}
	\begin{exampleblock}{\javapackageDrei{MP3Player}}
 		\begin{lstlisting}
public void playSongsRandomIntern() {
    Iterator<Music> iterator = musicLibrary.
            createRandomIterator();
    iterator.iterate(mp3 -> System.out.println(mp3));
}

public void playSongsByQueenIntern() {
    Iterator<Music> iterator = musicLibrary.
            createArtistIterator("Queen");
    iterator.iterate(mp3 -> System.out.println(mp3));
}
 		\end{lstlisting}
 	\end{exampleblock}
\end{frame}

\note{
	Der Vorteil eines internen Iterators ist, dass der Klient weniger über die Logik wissen muss.
	Dafür sind externe Iteratoren flexibler einsetzbar, weil der Klient die volle Kontrolle
	über den Iterationsprozess hat. Welche Variante die bessere ist, hängt vom Anwendungsfall ab. 
}

\frameWithImageOnly{\BILDERPFAD/iterator8.jpg}
\note{
	Was machen wir jetzt mit Änderungen, die während des Iterierens durchgeührt werden sollen? Bei 
	selbst geschrieben Iteratoren muss man sich darüber natürlich Gedanken machen, wenn man so
	etwas erlauben möchte. Wenn man es verbieten möchte, geht dies natürlich nur mit einer 
	ordentlichen Kapselung. Auf keinen Fall sollte man das Problem ignorieren, denn es könnte zu
	Laufzeitfehlern führen, die im Extremfall logische Fehler zur Folge haben. Im einfachsten Fall 
	werden die Elemente einfach im Iterator kopiert und Änderungen werden nur an dieser internen
	Liste durchgeführt. Das scheidet aus Performancegründen bei vielen Datensätzen aus. Man könnte -- wie 
	bei der ArrayList -- die entsprechenden Methoden im Iterator unterbringen. Dann müssen Iterator und
	Aggregat allerdings miteinander sprechen, denn es könnte ja sein, dass verschiedene Klienten
	zur gleichen Zeit über die Daten laufen und Änderungen vornehmen.  
}

\frameWithImageOnly{\BILDERPFAD/iterator9.jpg}
\note{
	Mit Iteratoren kann man auch durch komplexere Datenstrukturen laufen wie bspw. Bäume, die mithilfe
	eines Kompositummusters erstellt wurden. Dann muss sich der Iterator nicht nur die Position des
	aktuellen Elements merken, sondern auch den Pfad dorthin. Man könnte einen solchen Iterator auch
	rekursiv lösen. Dann müsste er sich nicht den gesamten Pfad merken, sondern der Pfad ergibt sich
	aus dem Aufrufstack. 
}